% gwebman.tex -- the GWEB manual.  Process with a plain-TeX engine:
%   pdftex gwebman.tex
% This is the GWEB analogue of CWEB's cwebman.tex.  It is a self-contained plain
% TeX document and does not need gwebmac.tex.

\hsize=6.3in \vsize=8.7in
\parindent=18pt
\tolerance=2000 \hbadness=2000

\font\titlefont=cmbx12 scaled\magstep2
\font\chapfont=cmbx12 scaled\magstep1
\font\secfont=cmbx12
\font\eightrm=cmr8

% --- inline code word: @ < > = etc. are safe in \tt; avoid braces inline ---
\def\cw#1{{\tt #1}}
\def\bs{\char92 }                             % a literal backslash, in \tt
\def\rn#1{$\langle\,${\it #1\/}$\,\rangle$}   % a refinement name, as in print

% --- structural macros ---
\newcount\chapno
\def\chapter#1{\par\penalty-300\vfil\eject\global\advance\chapno by1
  \leftline{\chapfont \the\chapno.\enspace #1}\nobreak\bigskip\noindent}
\def\section#1{\par\penalty-150\medskip\leftline{\secfont #1}\nobreak
  \smallskip\noindent}

% --- a verbatim code display (TeXbook recipe): | is the escape inside it ---
\def\begincode{\par\medskip\begingroup
  \def\do##1{\catcode`##1=12 }\dospecials
  \catcode`\|=0
  \obeylines \obeyspaces \tt \parindent=0pt \leftskip=18pt\relax}
\def\endcode{\endgroup\medskip\noindent}

\def\bull{\smallbreak\noindent\llap{$\bullet$\enspace}}

% =============================================================================
\centerline{\titlefont GWEB}
\medskip
\centerline{\secfont A Literate Programming System for Go}
\bigskip
\centerline{(modeled on Knuth and Levy's CWEB)}
\vskip 24pt
\noindent
This manual describes {\bf GWEB}, a system that lets you write a Go program and
its documentation as a single source file in the style of literate programming.
GWEB follows Donald Knuth and Silvio Levy's {\bf CWEB} as closely as possible;
to a first approximation it is ``CWEB with C replaced by Go.'' If you already
know CWEB, you can read this manual quickly, noting only the Go-specific points
in the final chapter.

\chapter{What literate programming is}
A {\it literate program\/} is written for human beings first and computers
second. Instead of arranging code in the order the compiler demands, you arrange
it in the order that best explains the ideas, breaking the program into small
labeled pieces and weaving prose between them. Two programs then process the
single source file in opposite directions:

\bull {\bf \cw{gtangle}} produces a Go source file for the compiler. It discards
all the commentary and reassembles the labeled pieces into the order Go requires.

\bull {\bf \cw{gweave}} produces a \TeX\ file for the reader. It typesets the
prose, pretty-prints the Go code (reserved words in {\bf bold}, identifiers in
{\it italic\/}), and builds an index and a list of code pieces with full
cross-references.

\noindent A GWEB source file conventionally has the extension \cw{.w}. From
\cw{foo.w} you obtain \cw{foo.go} with \cw{gtangle} and \cw{foo.tex} (then
\cw{foo.pdf}) with \cw{gweave}.

\chapter{The structure of a web}
A web is a sequence of {\it sections}. Material before the first section is
called {\it limbo}: it is plain \TeX\ that \cw{gweave} copies to the output
(typically \cw{\bs input gwebmac} and a title), and that \cw{gtangle} ignores.

Each section has up to three parts, always in this order:

\bull the {\it \TeX\ part}: commentary, written in ordinary \TeX;

\bull the {\it definition part}: zero or more \cw{@d}, \cw{@f}, or \cw{@s}
directives;

\bull the {\it code part}: the Go code of the section.

\noindent A section is opened by one of two control codes. \cw{@\ } (the
character \cw{@} followed by a space, tab, or newline) opens an ordinary
section; \cw{@*} opens a {\it starred\/} section, which begins a major group and
whose title --- the text up to the first period --- appears as a heading and in
the table of contents. Write \cw{@*1}, \cw{@*2}, $\ldots$ for subordinate
groups; \cw{@*} alone means depth~0.

The code part is opened either by \cw{@c}, which marks {\it unnamed\/} code
belonging to the program as a whole, or by a section name of the form
\cw{@<...@>=}, which {\it defines\/} a named section (see the next chapter).
Sections numbered automatically, starting at~1.

\section{A first example}
Here is a complete one-section web:
\begincode
@* Hello. Our first GWEB program.
@c
package main

import "fmt"

func main() {
    fmt.Println("hello, world")
}
|endcode
The starred section's title is ``Hello''. Running \cw{gtangle} on this file
yields a \cw{hello.go} that compiles and runs; running \cw{gweave} yields a
typeset document with ``1. Hello.'' as its heading.

\chapter{Named sections (refinements)}
The power of literate programming comes from {\it named sections}, which CWEB
and GWEB also call {\it refinements}. You may use a refinement before you define
it, and you may define it in several pieces scattered through the document.

A reference \cw{@<Name of section@>} stands for code defined elsewhere. A
definition \cw{@<Name of section@>=} introduces that code. When \cw{gtangle}
meets a reference it substitutes the corresponding definition (recursively); when
\cw{gweave} meets one it prints \rn{Name of section} followed by the number of
the defining section, so the reader can navigate.

\begincode
@ The skeleton names a refinement and fills it in below.
@c
package main

func main() {
    @<Greet the user@>
}

@ Here is the refinement itself.
@<Greet the user@>=
println("hi!")
|endcode

Several definitions with the same name are concatenated in the order they
appear, so you can grow a section incrementally. \cw{gweave} marks the second and
later pieces with ``$\mathrel+\equiv$'' and notes ``This code is also defined
in'' the other sections. Under the first definition it lists every section that
{\it uses\/} the refinement.

A name may be abbreviated by a trailing \cw{...}: \cw{@<Greet...@>} matches the
unique full name beginning ``Greet''.

\chapter{Producing more than one file}
By default \cw{gtangle} writes the unnamed sections to \cw{basename.go}. To
direct code to an additional file, define a section whose name is that file:
\begincode
@(go.mod@>=
module hello

go 1.22
|endcode
Each such file is written verbatim (after refinement expansion) to the named
path. This is handy for emitting a \cw{go.mod}, a second package file, or
generated data alongside the main program.

You can also split a large web across several source files. The directive
\cw{@i filename} includes another \cw{.w} file at that point, exactly as if its
text appeared inline.

\chapter{Change files}
Sometimes you must adapt a web you do not want to edit --- someone else's
program, or a version tailored to a particular system. For this both tools take
an optional second argument, a {\it change file}, exactly as CWEB does:
\begincode
gtangle foo.w foo.ch
gweave  foo.w foo.ch
|endcode
A change file is a list of changes, each of which names some lines to find in
the master source and the lines to put in their place:
\begincode
@x
  the lines to find
@y
  the lines to substitute
@z
|endcode
Anything outside an \cw{@x}\thinspace$\ldots$\thinspace\cw{@z} group is ignored,
so you can annotate the change file freely. The three controls must begin a
line. Changes are applied to the master (after \cw{@i} includes are expanded) in
the order they occur: at the first line that equals a change's first line, the
whole block must match --- trailing spaces aside --- and is then replaced. It is
an error if a change never matches, or matches its first line but not the rest,
so a change is anchored by quoting enough surrounding context to be unique.

Because changes are applied before anything else, they affect \cw{gtangle} and
\cw{gweave} identically: the program you build and the document you read stay in
step. The bundled \cw{wc.ch} turns the \cw{wc} example's output into CSV.

\chapter{The woven document}
\cw{gweave} writes a \TeX\ file that expects the macros in \cw{gwebmac.tex} (the
analogue of \cw{cwebmac.tex}); your limbo should therefore start with
\cw{\bs input gwebmac}. Make sure \TeX\ can find that file, for example with
\hbox{\cw{TEXINPUTS=/path/to/gweb/tex: pdftex foo.tex}}.

The woven output prints each section in source order, pretty-printing the code:
reserved words and predeclared type names are set in bold, other identifiers in
italic, string and rune literals in typewriter, and comments in roman. The
original indentation of the Go code is preserved.

\section{Inline code, indexing, and formatting}
Within commentary you may insert a short piece of Go code by surrounding it with
vertical bars: writing \hbox{\cw{\char124 x := f()\char124}} typesets it like
code. (Use \cw{\bs\char124} for a literal bar.)

Three control codes add hand-written entries to the index, each ending with
\cw{@>}:

\bull \cw{@\char94 text@>} --- typeset the entry in roman;

\bull \cw{@.text@>} --- typeset it in typewriter;

\bull \cw{@:text@>} --- typeset it with the user macro \cw{\bs 9}.

\noindent Finally, \cw{@f a b} tells \cw{gweave} to format identifier \cw{a} the
way it formats \cw{b} (for instance to make a type name bold); \cw{@s} does the
same but omits the entry from the index.

\chapter{Running the programs}
Build the tools once with \cw{make build} (or \cw{go install ./...}), then:
\begincode
gtangle foo.w        # -> foo.go
gweave  foo.w        # -> foo.tex
pdftex  foo.tex      # -> foo.pdf
|endcode
Both tools accept \cw{-o dir} to choose the output directory. \cw{gtangle} runs
\cw{gofmt} on its result, so a correct web always tangles to canonically
formatted Go; if the assembled program is not valid Go, \cw{gtangle} keeps the
raw output and prints a warning.

Both tools also print diagnostics, each tagged with a \cw{file:line} location:
an unterminated control code, a reference to an undefined section, an ambiguous
\cw{...} abbreviation, or a section that is defined but never used. These are
warnings and do not stop processing --- although \cw{gtangle} will still fail if
it actually has to expand a reference that was never defined.

\chapter{Differences from CWEB and Go-specific notes}
GWEB keeps CWEB's control codes and behavior wherever they make sense. The main
differences follow from the move from C to Go.

\bull {\bf No preprocessor.} Go has no macros, so \cw{@d} has no analogue: it is
accepted for compatibility but ignored. Prefer Go \cw{const} and \cw{func}
declarations. The formatting directives \cw{@f} and \cw{@s} remain useful: for
example \cw{@f Counts int} sets the user type \cw{Counts} in bold, the way
\cw{gweave} sets predeclared types.

\bull {\bf Formatting from \cw{gofmt}.} Rather than reformatting code itself,
\cw{gtangle} delegates to \cw{gofmt}; the tangled file is always tidy.

\bull {\bf Layout in the woven output.} \cw{gweave} mirrors the source's own
spacing and indentation instead of re-deriving layout from a full grammar as
\cw{cweave} does for~C. Because gofmt-formatted Go already encodes the grammar in
its spacing, this reproduces gofmt exactly --- pointer \cw{*T} versus the product
\cw{a * b}, the slice \cw{[]T} versus the index \cw{a[i]} --- with no parsing.
Long code lines wrap at the spaces between tokens, with continuation lines
indented. Write the code in your sections in gofmt style for the best results.

\bull {\bf Heuristic definitions.} An identifier is marked as ``defined'' in the
index (its section number underlined) when it follows \cw{func}, \cw{var},
\cw{const}, or \cw{type}, or stands just left of \cw{:=}. This is a heuristic,
not a full type analysis.

\vfill
\centerline{\eightrm GWEB --- a literate programming system for Go.}
\bye
